\scanpage{7}
\mspage{2}
$\tau$, dont l'opération sur $\pi_{0,3}$ peut se décrire par
\[
\tau(l_0)=l_0^{-1},\quad\tau(l_1)=l_1^{-1},\quad\tau(l_\infty)=l_1l_0=\operatorname{int}(l_1)l_\infty^{-1}
\]
C'est donc un autom du groupe à lacets, de signature $-1$, satisfaisant
\[
\tau^2=1
\]
On vérifie qu'il commute à l'action de $\unsure{\mathfrak S_3|\mathcal S_3}$ mod. autom extérieurs --- ce qui donne une op. extérieure de $\unsure{\mathfrak S_3|\mathcal S_3}\times\{\pm1\}=\Gamma_{0,3}$, qui n'est autre que l'action canonique.

On s'intéresse maintenant aux autom extérieurs qui commutent à $\rho$. Notons que tout autom extérieur \inserted{qui normalise la classe de conj. de $l_0$} provient d'un autom $u$ qui normalise \deleted{\unsure{un}} $L_0$ ($\simeq\hat{\mathbb Z}$ engendré par $l_0$) et $u$ est unique mod composition par un $\lambda\in L_0$. Si $p$ \inserted{$=\chi(u)$}, $\in\hat{\mathbb Z}^{*}$ \deleted{\unsure{est}} est le multiplicateur de $u$, on aura donc
\[
(*)\qquad u(l_0)=l_0^p
\]
Soit $g\in\pi$ tel que :
\[
(**)\qquad u\rho=\deleted{\rho u}\operatorname{int}(g)\rho u,
\]
ce qui détermine $g$ de façon unique, $u$ étant
\mspage{6}
